\documentclass[12pt]{article}
\usepackage{amsmath,amssymb,amsthm}
\usepackage[margin=1in]{geometry}
\usepackage{booktabs}
\usepackage{hyperref}

\newtheorem{theorem}{Theorem}[section]
\newtheorem{corollary}[theorem]{Corollary}
\newtheorem{proposition}[theorem]{Proposition}
\newtheorem{definition}[theorem]{Definition}
\theoremstyle{remark}
\newtheorem{remark}[theorem]{Remark}

\newcommand{\CC}{\mathbb{C}}
\newcommand{\RR}{\mathbb{R}}
\newcommand{\QQ}{\mathbb{Q}}
\newcommand{\nS}{n_{\mathrm{S}}}
\newcommand{\nT}{n_{\mathrm{T}}}
\newcommand{\Lean}{\textsc{Lean\,4}}

\title{The Birch and Swinnerton-Dyer Conjecture\\
  as the $(3,2)$ Correspondence:\\
  Taniyama--Shimura from Simplicial Geometry}
\author{Mingu Jeong\\Independent Researcher
\and Claude\\Anthropic}
\date{\today}

\begin{document}
\maketitle

\begin{abstract}
The Taniyama--Shimura theorem (modularity) establishes a
correspondence between elliptic curves (degree 3, genus 1)
and weight-2 modular forms (SL$(2, \ZZ)$ action). We identify
this as the $(3,2)$ decomposition of $\CC^5$: the elliptic
curve lives in $\CC^3$ (spatial sector, $\nS = 3$), the
modular form in $\CC^2$ (temporal sector, $\nT = 2$), and
the shared L-function is $\mathrm{ref} \circ \mathrm{incl}
= G_{ij}$ (the Gram matrix). The genus formula
$(n{-}1)(n{-}2)/2$ equals the CKM CP-phase count,
giving genus 1 and 1 CP phase at $n = 3$ simultaneously.
Both the BSD conjecture and CP violation arise from the
same ``3'': the minimum cycle length for a gauge-invariant
phase (Bargmann invariant). Spectral complexity:
standard $(1,4)$, physical $(1,2)$ \cite{Paper9}.
\end{abstract}

\section{Introduction}

The BSD conjecture relates the rank of an elliptic curve
$E/\QQ$ to the order of vanishing of $L(E, s)$ at $s = 1$.
This requires the modularity theorem (Taniyama--Shimura),
which states $L(E, s) = L(f, s)$ for a weight-2 modular form $f$.

We show this correspondence has a structural origin:
the $(3,2)$ decomposition of $\CC^5$.

\section{Taniyama--Shimura as $(3,2)$}

\begin{center}
\begin{tabular}{lll}
\toprule
Taniyama--Shimura & & DRLT \\
\midrule
Elliptic curve (degree 3) & $\leftrightarrow$
  & $\CC^3$ sector (spatial, $\nS = 3$) \\
Modular form (SL$(2)$, weight 2) & $\leftrightarrow$
  & $\CC^2$ sector (temporal, $\nT = 2$) \\
$L(E,s) = L(f,s)$ & $\leftrightarrow$
  & $\mathrm{ref} \circ \mathrm{incl} = G_{ij}$ \\
GL$_2$ representation (dim 2) & $\leftrightarrow$
  & $\dim_\RR(\CC) = 2 = \nT$ \\
\bottomrule
\end{tabular}
\end{center}

\section{Genus = CP Phases}

\begin{theorem}
The genus of a degree-$n$ smooth projective curve and the
number of CP phases in an $n \times n$ CKM matrix are given
by the \emph{same formula}:
\[
  g(n) = \frac{(n-1)(n-2)}{2}
  = \text{\# CP phases in $n$-generation CKM}.
\]
\end{theorem}

\begin{center}
\begin{tabular}{cccc}
\toprule
$n$ & $(n{-}1)(n{-}2)/2$ & Fermat curve & CKM matrix \\
\midrule
2 & 0 & genus 0 ($\infty$ Pythagorean) & 0 CP phases \\
3 & 1 & genus 1 (elliptic, FLT) & 1 CP phase \\
4 & 3 & genus 3 (Faltings) & 3 CP phases \\
\bottomrule
\end{tabular}
\end{center}

The transition at $n = 3$ is the \emph{same} for both:
the minimum cycle (Bargmann invariant) requires 3 elements.

\section{The Bargmann Bridge}

For $k$ vectors $\psi_1, \ldots, \psi_k$ in $\CC^d$:
\[
  B_k = \langle\psi_1|\psi_2\rangle
        \langle\psi_2|\psi_3\rangle
        \cdots \langle\psi_k|\psi_1\rangle.
\]
Under rephasing $\psi_j \to e^{i\alpha_j}\psi_j$:
phases telescope, so $\arg(B_k)$ is gauge-invariant.

$k = 2$: $B_2 = |\langle\psi_1|\psi_2\rangle|^2 \in \RR_+$
(no phase). $k = 3$: $B_3 \in \CC$ (has phase = CP violation).

\begin{corollary}
3 = minimum cycle = minimum for gauge-invariant phase
= minimum genus = minimum CP.
MSUA's ``3 layers for meaning'' = CKM's ``3 for matter.''
\end{corollary}

\section{Spectral Complexity}

BSD standard: $(h, l) = (1, 4)$ (rank at $s = 1$ needs $\infty$).
BSD physical: $(h, l) = (1, 2)$ (the $(3,2)$ structure is
combinatorial, Level 2) \cite{Paper9}.

\section{Machine Verification}

\texttt{BSDPoincare.lean}: \texttt{bsd\_is\_32},
\texttt{genus\_eq\_cp}, \texttt{cp\_transition},
\texttt{seven\_millennium}. Zero \texttt{sorry}.

\begin{thebibliography}{9}
\bibitem{Paper9} M.~Jeong and Claude, ``Spectral Complexity,''
  DRLT Paper~9 (2026).
\bibitem{Wiles} A.~Wiles, ``Modular elliptic curves and
  Fermat's Last Theorem,'' \emph{Ann.\ Math.} (1995).
\end{thebibliography}
\end{document}